%
% 0_Einleitung.tex -- Einleitung zum Kapitel Resultante
%
% (c) 2026 Jakob Gierer, OST Ostschweizer Fachhochschule
%
% !TEX root = ../../buch.tex
% !TEX encoding = UTF-8
%
\section{Einleitung\label{resultante:section:einleitung}}
\kopfrechts{Einleitung}

Für die Lösung von Integralen greifen Ingenieur:innen sowie Studierende oft auf Computeralgebrasysteme (CAS) zurück.
\index{Computeralgebrasystem}%
Damit CAS-fähige Rechner Integrale lösen können, werden effiziente Algorithmen benötigt.
Integranden werden zunächst in einen polynomiellen Teil und einen rationalen Teil aufgeteilt.
Dieses Kapitel behandelt nur Methoden zur Bestimmung der Stammfunktion des rationalen Teils.
Dabei besteht der rationale Teil des Integrals aus echt gebrochenen Funktionen. 
Für die Berechnung dieser Integrale müssen die Nullstellen des Nennerpolynoms der Integranden bestimmt werden.

In Abschnitt \ref{buch:ratint:section:rothsteintrager} wird der Satz von Rothstein-Trager beschrieben, der für die Bestimmung der Nullstellen in Kapitel \ref{chapter:intalg} den euklidischen Algorithmus nutzt.
\index{Rothstein-Trager, Satz von}%
\index{Satz!von Rothstein-Trager}%
In diesem Kapitel soll eine weitere Methode zur Berechnung von Nullstellen vorgestellt werden, die Resultante.
In verschiedenen Lehrbüchern wird die Resultante als
\begin{equation}
    \Res(f,g)=\det(\syl(f,g))
    \label{resultante:equation:definition}
\end{equation}
definiert. Aber was ist \(\syl\)? 
Und was ist die Resultante?

Was genau die Resultante ist und welche nützlichen Eigenschaften sie hat, beschreibt dieses Kapitel.
Die Beschreibungen und Ideen in diesem Kapitel basieren auf \cite{resultante:Linalg} und \cite{resultante:SymInt}.
